% !TeX root = ../informe.tex
% ===========================================================================
%  Datos de la portada y metadatos del PDF.
%  Este es el único archivo que hay que editar con los datos del equipo.
% ===========================================================================

\newcommand{\tituloInforme}{Sistema de Reserva de Canchas Deportivas}
\newcommand{\subtituloInforme}{Arquitectura de Microfrontends y Microservicios}
\newcommand{\nombreCurso}{Desarrollo de Aplicaciones Empresariales}
\newcommand{\nombrePrograma}{Maestría en Ingeniería de Software}
\newcommand{\nombreUniversidad}{Universidad Politécnica Salesiana}
\newcommand{\tipoTrabajo}{Proyecto Integrador}

% --- Integrantes del equipo ------------------------------------------------
% Un par nombre/correo por integrante. Dejar vacíos los que sobren: la
% portada omite automáticamente los que estén en blanco.
\newcommand{\autorUno}{Pablo Jara Muñoz}
\newcommand{\correoUno}{pablojaramunoz@gmail.com}

\newcommand{\autorDos}{Miguel Lara}
\newcommand{\correoDos}{miguel.larasj@gmail.com}

\newcommand{\autorTres}{Sebastián Uyaguari}
\newcommand{\correoTres}{sebastianuyaguari@gmail.com}

\newcommand{\autorCuatro}{Adrián Espinoza}
\newcommand{\correoCuatro}{adrian.espinoza1292@gmail.com}

\newcommand{\nombreDocente}{Christian Merchán Millán}
\newcommand{\ciudadFecha}{Ciudad, \today}

% --- Metadatos del PDF -----------------------------------------------------
\hypersetup{
    pdftitle  = {\tituloInforme},
    pdfauthor = {\autorUno, \autorDos, \autorTres, \autorCuatro},
    pdfsubject= {\nombreCurso\ --- \nombrePrograma},
    pdfkeywords = {microservicios, microfrontends, Module Federation,
                   Spring Boot, PostgreSQL, Docker}
}
